\inicializarSeccion{Introducción}

\tab En 1949, el programa de Investigaciones Cooperativas Oceánicas Pesqueras de California(CalCOFI) fue creado al  notar el colapso de la pesquería de sardina en la costa de California. Este tuvo como objetivo estudiar y analizar las causas ambientales y biológicas relacionadas con esta problemática significativa. 

En la actualidad, el programa evalúa variables físicas, químicas y biológicas del océano para entender el efecto provocado por el cambio climático y actividades humanas en los ecosistemas marinos y los recursos pesqueros (\cite{calcofi_conference}). Debido a este enfoque, en la sustentabilidad marina, variables como temperatura, oxígeno disuelto, salinidad y nutrientes permiten evaluar el estado de salud de los ecosistemas oceánicos e identificar posibles condiciones de estrés ambiental.

Estudios recientes identifican que el incremento de la temperatura oceánica, las olas de calor y la disminución del oxígeno disuelto representan amenazas para la biodiversidad marina y la sostenibilidad de las zonas costeras (\cite{sanderson_2024_missing_values}).  Además, reportes del 2024 de  CalCOFI indican que los ecosistemas marinos podrían experimentar temperaturas más cálidas y menores concentraciones de oxígeno disuelto, generando alteraciones en la productividad biológica y en la distribución de especies marinas.

\textbf{Objetivo de indagación:} Analizar cómo varían el oxígeno disuelto, la saturación de oxígeno, la temperatura, la salinidad y la concentración de nutrientes en función de la profundidad y el fosfato, con el propósito de identificar patrones ambientales asociados a posibles condiciones de estrés ecológico y cambios en la productividad biológica marina